\chapter{Security Evaluation}\label{ch:security}

This chapter evaluates the system against the threat model of Chapter~\ref{ch:req} and
reports two results that the project did not set out to find: that the layered configuration
fails under one specific compromise, and that the mechanism defending against replay can lock out
the legitimate owner. Both are reported as findings about this design rather than as
limitations of the evaluation.

\section{The threat model, restated with outcomes}

Each scenario is defined by \emph{what the adversary holds}, not by what mechanism is deployed.
The access column is material: it separates a scenario requiring physical possession of the
device from one that can be mounted from anywhere on the network.

\begin{longtable}{@{}p{0.05\textwidth}p{0.26\textwidth}p{0.13\textwidth}p{0.42\textwidth}@{}}
\toprule
\textbf{ID} & \textbf{Adversary holds} & \textbf{Access} & \textbf{Outcome against the layered
build} \\
\midrule
\endhead
A & nothing                     & remote   & Refused at the broker; the connection is not
established \\
B & the one shared credential   & remote   & \textbf{Defeats} a shared-credential deployment;
refused by per-client authorisation \\
C & controller broker credential& remote   & Delivered; refused by the message authentication
code \\
D & device broker credential    & remote   & Refused by the access control list; the device
identity may not publish commands \\
E & device credential + command key & \textbf{physical} & \textbf{Survived}: refused by the
access control list despite a valid signature \\
E$'$ & controller credential + command key & \textbf{remote} & \textbf{Defeats the layered
build}: neither layer refuses it \\
\bottomrule
\end{longtable}

\subsection{Why E and E$'$ are the decisive pair}\label{sec:sec-eprime}

Scenarios E and E$'$ carry \emph{the same cryptographically valid signature}. They differ only in
the broker identity that transmits it.

Under E the signature arrives from the device identity. The access control list grants that
identity \code{read} on \code{home/alarm/+/cmd} and \code{write} on \code{home/alarm/+/stat};
it does not grant \code{write} on the command topic. The broker refuses the publication, so the
valid signature never reaches the device.

Under E$'$ the same signature arrives from the controller identity, which must hold
\code{write} on the command topic for the system to function at all. The broker delivers it, and
the message authentication code verifies --- correctly, because the signature is valid. The
command executes.

\textbf{What stops scenario E is therefore the access control list and not the signature.} This
cannot be established by reasoning about the cryptography, and it is the result that justified
running experiment E6 against real broker software rather than modelling broker behaviour. Both
scenarios were executed against Mosquitto 2.0.18 over TLS 1.2 with a private certificate
authority; the verdicts in Section~\ref{sec:res-broker} are the broker's own.

\subsection{An asymmetry in cost}

E requires physical possession of the device and a flash dump. E$'$ requires compromise of the
controller host --- a general-purpose machine, typically a Raspberry Pi or a server, reachable
over the network and administered by software with its own vulnerability history.

The layered configuration therefore \textbf{survives the expensive attack and falls to the cheap
one.} That is the opposite of the ordering a defender would choose, and it is stated plainly
because an earlier draft of this work did not state it.

\section{Flash extraction and what is actually protected}\label{sec:sec-flashdump}

\subsection{The threat}

The ESP32's SPI flash can be read over the serial interface with standard vendor tooling unless
flash encryption and secure boot are enabled. This firmware does not enable them. Anyone with
physical possession of the device can therefore recover the entire image, which contains both the
device's broker credential and \code{SECRET\_CMD\_KEY}.

\subsection{Why the device survives it anyway}

This is a genuine positive result and is worth stating precisely, because it is easy to
overstate.

The device survives a flash dump \emph{not} because the key is protected --- it is not --- but
because \textbf{the identity whose secrets were extracted is not authorised to issue commands.}
The access control list is deny-by-default and grants the device identity publication rights only
on the status topic. An adversary holding everything in the device's flash can sign a
valid \code{DISARM} command and still cannot get it accepted, because the broker refuses to
carry it.

This is a defence-in-depth result in the strict sense: a complete failure of secret
confidentiality at one layer is contained by authorisation at another.

\subsection{What it does not mean}

It does not mean key extraction is harmless. It means key extraction \emph{from the device} is
contained. Key extraction from the controller host is scenario E$'$, which is not contained. The
protection comes from \emph{which} identity holds the key, not from the key being hard to obtain.

Three consequences follow and are recorded in Chapter~\ref{ch:limitations}: a deployment that
issues one credential to both parties collapses E$'$ and E into the same scenario and loses the
protection entirely; a deployment whose access control list grants the device identity write
access on the command topic, for convenience during commissioning, loses it as well; and flash
encryption with secure boot would address the extraction itself rather than only containing it.

\section{The compile-time guard}\label{sec:sec-guard}

\subsection{The original defect}

The firmware originally admitted remote disarm when the build had transport security and a
broker login, on the reasoning that a client which has authenticated to the broker over TLS is
trustworthy. That reasoning is wrong for a specific and documented reason: in MQTT 3.1.1 the
broker \emph{re-originates} messages. The subscriber receives a message from the broker, not from
the publisher, and the protocol carries no publisher identity to the subscriber. Transport
authentication tells the broker who connected; it tells the device nothing about who authored a
command.

Broker credentials are bearer tokens. One leaked credential makes the adversary a legitimate
client, and the ablation measured the consequence: alarm activation falls from 300/300 to 50/300,
with the adversary disarming the device in 246 of 300 trials
(Section~\ref{sec:res-ablation}).

\subsection{The repair, and why it is a build-time error}

The defective configuration is now rejected at compile time:

\begin{lstlisting}[style=cpp,caption={The build guard. The configuration cannot be built
accidentally.},label={lst:sec-guard}]
#ifndef ALLOW_REMOTE_DISARM
#  define ALLOW_REMOTE_DISARM 0
#endif
#if ALLOW_REMOTE_DISARM && !AUTH_COMMANDS && !defined(ACKNOWLEDGE_UNAUTHENTICATED_DISARM)
#  error "ALLOW_REMOTE_DISARM without AUTH_COMMANDS 1 permits anyone who can publish to the \
command topic to switch the alarm off, and TLS plus a broker login does not prevent it (see \
docs/SECURITY_ANALYSIS.md). Set AUTH_COMMANDS 1, or define \
ACKNOWLEDGE_UNAUTHENTICATED_DISARM to accept that risk deliberately."
#endif
\end{lstlisting}

Three design decisions in that guard are deliberate.

\textbf{It is an error, not a warning.} A warning in an Arduino build is lost in several hundred
lines of framework output.

\textbf{The escape hatch is named for what it concedes.} The macro is
\code{ACKNOWLEDGE\_UNAUTHENTICATED\_DISARM}, not \code{FORCE} or \code{OVERRIDE}. A builder who
defines it has written the concession into their build command.

\textbf{The escape hatch exists.} The defective configuration must remain buildable, because it
is the research baseline. Deleting it would make the measured comparison irreproducible. This is
the one place in the firmware where a known-insecure configuration is retained on purpose, and
the reason is recorded next to it.

\subsection{A second guard: the published key}

A related defect was found and guarded during the same review. Authentication that verifies
against a published key is not authentication. The repository's \code{SECRET\_CMD\_KEY} default
is a public demonstration string, so a build enabling \code{AUTH\_COMMANDS} without overriding it
would log every command as authenticated while any reader of the repository could forge one. That
configuration is now also rejected at build time.

Three negative build tests in \code{tests/host/run\_tests.sh} require each defective
configuration to \emph{fail to compile}. A build that succeeds is a test failure.

\section{Availability as a security property}

The ablation treats operator availability as a measured security property rather than a
usability concern, and two arms show why.

\textbf{The rate limiter.} A token bucket on the command path preserves the safety function and
leaves replay resistance unchanged, and costs \textbf{41\% of legitimate operator commands}
(737 of 1,250 admitted, Fisher exact $p = 3.5 \times 10^{-184}$). An alarm that refuses two in
five owner commands is not a secure alarm; it is a denial of service the defender installed.

\textbf{Owner lockout under key compromise.} Two configurations admit \textbf{0 of 1,250}
operator commands. The adversary signs commands with a counter far above the operator's; the
device persists that counter, as replay resistance requires; every subsequent legitimate command
then carries a lower counter and is correctly rejected as a replay.

The mechanism behaved exactly as specified. The specification was incomplete: it did not
consider that an adversary who can advance the counter can also exhaust the operator's ability
to use it. This was not predicted by this project, though the mechanism itself is not new: Wara
and Yu~\cite{wara2020} state it as background in the ZigBee setting, without a citation and
without an experiment. What is contributed here is the measurement --- a complete rather than
partial lockout, 0 of 1,250 --- and the resulting observation that replay resistance and owner
availability are coupled under key compromise. It is reported as a negative result about a design
decision made in this project, and a candidate repair --- a counter-resynchronisation path
authenticated by the local button, which is physically present and outside the adversary's reach
--- is identified in Chapter~\ref{ch:limitations} as future work. It has not been implemented or
evaluated.

\section{Properties established, and by which instrument}

\begin{longtable}{@{}p{0.31\textwidth}p{0.19\textwidth}p{0.42\textwidth}@{}}
\toprule
\textbf{Property} & \textbf{Evidence} & \textbf{Scope of the claim} \\
\midrule
\endhead
No unauthorised disarm with command authentication & 6,000 trials + exhaustive model checking
& Within the harness's adversary model and the model's bounds \\
Replay resistance                  & 3,600 trials, 0 successes & Against capture-and-replay of
observed ciphertext \\
Memory safety of the parser        & 3,000,000 fuzz iterations under both sanitizers & Host
build; not the Xtensa target \\
Flash dump contained               & live broker, scenario E & Contingent on per-client
authorisation and on the device identity lacking command write access \\
Logical isolation of the safety task & exhaustive model checking, 660 states & Logical only;
\textbf{no physical timing claim} \\
Specification conformance          & 19,529 transitions, 0 violations & Three ablation modes,
with a negative control \\
\bottomrule
\end{longtable}

\section{Properties explicitly not established}

\begin{itemize}[leftmargin=1.6em,itemsep=4pt]
  \item \textbf{Physical timing invariance.} Not measured. Hardware conditions A--D were not
        executed (Appendix~\ref{app:hardware}). Every isolation claim in this work is logical.
  \item \textbf{On-device memory safety.} The fuzzing campaign ran against a host build with
        AddressSanitizer and UndefinedBehaviorSanitizer, neither of which exists for the Xtensa
        target. Memory safety of the compiled Xtensa binary is not established.
  \item \textbf{Resistance to radio-layer attack.} No RF-shielded enclosure was available, so no
        deauthentication or jamming experiment was performed against real hardware.
  \item \textbf{Resistance to side-channel or fault-injection attack.} Out of scope; no
        equipment.
  \item \textbf{Prevalence of the guard defect in deployed systems.} The corpus study returned
        zero hits for the precise rule and nine confirmed cases for the looser one. That is an
        existence claim over a biased 135-project sample, not a prevalence estimate
        (Section~\ref{sec:res-corpus}).
  \item \textbf{Security of the controller host.} Treated as trusted infrastructure throughout.
        Scenario E$'$ is precisely what happens when that assumption fails, and the layered build
        does not survive it.
\end{itemize}
